\documentclass{article}

\usepackage[spanish]{babel}
\usepackage[utf8]{inputenc}
\usepackage[T1]{fontenc}
\usepackage{tikz}
\usetikzlibrary{positioning}
\usepackage{multirow}
\usepackage[margin=1cm, landscape]{geometry}

\begin{document}

\begin{figure}[h]
\centering
\begin{tikzpicture}[node distance=2cm]

%======================= NODOS (TABLAS) =======================

% Usuario (Centro)
\node (Usuario) [inner sep=0pt] {
\begin{tabular}{|l|l|}
\hline
\multicolumn{2}{|c|}{\textbf{Usuario}}\\\hline
\textbf{PK} & id\\\hline
           & nombre\\
           & email\\
           & password\\
           & rol\\
           & fecha\_registro\\\hline
\end{tabular}
};

% Vendedor (Izquierda de Usuario)
\node (Vendedor) [left=3cm of Usuario, inner sep=0pt] {
\begin{tabular}{|l|l|}
\hline
\multicolumn{2}{|c|}{\textbf{Vendedor}}\\\hline
\textbf{PK} & id\\\hline
\textbf{FK} & id\_usuario\\\hline
            & nombre\_tienda\\
            & rating\\
            & ubicacion\\
            & estado\\\hline
\end{tabular}
};

% Pedido (Derecha de Usuario)
\node (Pedido) [right=3cm of Usuario, inner sep=0pt] {
\begin{tabular}{|l|l|}
\hline
\multicolumn{2}{|c|}{\textbf{Pedido}}\\\hline
\textbf{PK} & id\\\hline
\textbf{FK} & id\_usuario\\\hline
            & total\\
            & estado\\
            & fecha\\\hline
\end{tabular}
};

% Producto (Abajo de Vendedor)
\node (Producto) [below=3cm of Vendedor, inner sep=0pt] {
\begin{tabular}{|l|l|}
\hline
\multicolumn{2}{|c|}{\textbf{Producto}}\\\hline
\textbf{PK} & id\\\hline
\textbf{FK} & id\_vendedor\\
\textbf{FK} & id\_categoria\\\hline
            & titulo\\
            & precio\\
            & stock\\
            & marca\\
            & descripcion\\\hline
\end{tabular}
};

% Reseña (Abajo de Usuario)
\node (Resena) [below=3cm of Usuario, inner sep=0pt] {
\begin{tabular}{|l|l|}
\hline
\multicolumn{2}{|c|}{\textbf{Reseña}}\\\hline
\textbf{PK} & id\\\hline
\textbf{FK} & id\_usuario\\
\textbf{FK} & id\_producto\\\hline
            & calificacion\\
            & comentario\\
            & fecha\\\hline
\end{tabular}
};

% DetallePedido (Abajo de Pedido)
\node (DetallePedido) [below=3cm of Pedido, inner sep=0pt] {
\begin{tabular}{|l|l|}
\hline
\multicolumn{2}{|c|}{\textbf{DetallePedido}}\\\hline
\multirow{2}{*}{\textbf{PK FK}} & id\_pedido\\\cline{2-2}
                                & id\_producto\\\hline
                                & cantidad\\
                                & precio\_unitario\\\hline
\end{tabular}
};

% Categoría (Abajo de Producto)
\node (Categoria) [below=2cm of Producto, inner sep=0pt] {
\begin{tabular}{|l|l|}
\hline
\multicolumn{2}{|c|}{\textbf{Categoría}}\\\hline
\textbf{PK} & id\\\hline
            & nombre\\
            & icono\\\hline
\end{tabular}
};

%======================= RELACIONES =======================

% Usuario - Vendedor (1 : 1)
\draw (Usuario.west) -- node[above,pos=0.2]{1} node[above,pos=0.8]{1} (Vendedor.east);

% Usuario - Pedido (1 : N)
\draw (Usuario.east) -- node[above,pos=0.2]{1} node[above,pos=0.8]{N} (Pedido.west);

% Usuario - Reseña (1 : N)
\draw (Usuario.south) -- node[right,pos=0.2]{1} node[right,pos=0.8]{N} (Resena.north);

% Vendedor - Producto (1 : N)
\draw (Vendedor.south) -- node[right,pos=0.2]{1} node[right,pos=0.8]{N} (Producto.north);

% Pedido - DetallePedido (1 : N)
\draw (Pedido.south) -- node[right,pos=0.2]{1} node[right,pos=0.8]{N} (DetallePedido.north);

% Producto - DetallePedido (1 : N) - Cruzada larga
\draw (Producto.east) -- node[above,pos=0.1]{1} node[above,pos=0.9]{N} (DetallePedido.west);

% Producto - Reseña (1 : N)
\draw (Producto.east) -- node[above,pos=0.2]{1} node[above,pos=0.8]{N} (Resena.west);

% Categoría - Producto (1 : N)
\draw (Categoria.north) -- node[right,pos=0.2]{1} node[right,pos=0.8]{N} (Producto.south);

\end{tikzpicture}
\caption{Modelo Relacional - AppleAura}
\end{figure}

\end{document}
